\begin{vidu}%[3P3-15.3K][Loai1] 
	Điện năng được truyền tải đi xa với công suất không đổi $\mathcal{P}$ và hiệu điện thế hiệu dụng $U = 10\ kV$, hiệu suất của quá trình truyền tải $H = 80\%$. Để hiệu suất của quá trình truyền tải đạt được là $95\%$, thì phải tăng hiệu điện thế lên đến giá trị
	\choice
	{\True 20 kV}
	{25 kV}
	{30 kV}
	{25 kV}
	\loigiai{
		\begin{longtable}{|p{0.22\textwidth}|p{0.22\textwidth}|p{0.22\textwidth}|p{0.22\textwidth}|}
			\hline %========================
			P&	U	&$\Delta P$&$P_t$\\
			\hline %========================
			100&	10k&	20&	80\\
			\hline %========================
			100	& $U_2$	&5	&95\\
			\hline %========================
		\end{longtable}
		Công suất nơi phát không đổi: 
		$\Delta P\sim \dfrac{1}{U^2}\Rightarrow \dfrac{\Delta P_1}{\Delta P_2}=\dfrac{U_2^2}{U_1^2}\Leftrightarrow \dfrac{20}{5}=\dfrac{U_2^2}{10^2}\Leftrightarrow U_2=20\ kV$.
	}
\end{vidu}
\begin{vidu}%[3P3-15.3K][Loai1] 
	Điện năng được truyền tải từ nơi phát đến một khu dân cư bằng đường dây một pha với hiệu suất truyền tải là 90\%. Coi hao phí điện năng chỉ do tỏa nhiệt trên đường dây và không vượt quá 20\%. Nếu công suất sử dụng điện của khu dân cư này tăng 20\% và giữ nguyên điện áp ở nơi phát thì hiệu suất truyền tải điện năng trên chính đường dây đó là
	\choice
	{85,8\%}
	{\True 87,7\%}
	{89,2\%}
	{92,8\%}
	\loigiai{
		\begin{itemize}
			\item Điện áp nơi phát không đổi: $\Delta P\sim P^2\Rightarrow \dfrac{10}{a}=\left(\dfrac{100}{108+a}\right)^2\Leftrightarrow \begin{cases}
		a=768,8\\
		a=15,17
		\end{cases}$
		\item Để hao phí không vượt quá 20\% thì $a = 15,17$.
		\item Hiệu suất truyền tải lúc sau là: $H_2=\dfrac{108}{108+15,17}\approx 87,7\%$.
	\end{itemize}
	}
\end{vidu}
\begin{vidu}%[3P3-15.3K][Loai1] 
	Truyền tải điện năng đến nơi tiêu thụ bằng đường dây một pha. Hệ số công suất của mạch điện không đổi. Nếu điện áp ở trạm phát là $220\ V$ thì hiệu suất truyền tải điện năng là $80\% $. Để hiệu suất truyền tải là $92\%$ mà công suất truyền đến nơi tiêu thụ vẫn không thay đổi thì điện áp hiệu dụng ở trạm phát là
	\choice
	{\True $324\ V$}
	{$404\ V$}
	{$741\ V$}
	{$66\ V$}
	\loigiai{
		$\Delta P=\dfrac{P^2R}{U^2{\cos}^2\varphi};P'=P-\Delta P;H=\dfrac{P'}{P}$
		\begin{longtable}{|p{0.22\textwidth}|p{0.22\textwidth}|p{0.22\textwidth}|p{0.22\textwidth}|}
			\hline %========================
			$P$ &
			$U$ &
			$\Delta P$ 
			\\
			\hline %========================
			$100$ &
			$U_1=220$ &
			$20$ 
			\\
			\hline %========================
			$\dfrac{2000}{23}$ &
			$U_2$ &
			$\dfrac{160}{23}$ 
			\\
			\hline %========================
		\end{longtable}
		Vì $R,\cos \varphi $ không đổi nên $\dfrac{\Delta P_1}{\Delta P_2}=\dfrac{P_1^2}{P_2^2}.\dfrac{U_2^2}{U_1^2}\Rightarrow \dfrac{20}{\dfrac{160}{23}}=\dfrac{100^2}{{\left(\dfrac{2000}{23}\right)}^2}.\dfrac{U_2^2}{220^2}\Rightarrow U_2=324\ V$.
	}
\end{vidu}
\begin{vidu}%[3P3-15.4G][Loai2] 
	Một máy phát điện xoay chiều một pha có công suất phát điện và điện áp hiệu dụng ở hai cực máy phát đều không đổi. Điện năng được truyền từ máy phát đến nơi tiêu thụ trên đường dây có điện trở không đổi. Coi hệ số công suất của mạch này luôn bằng 1. Hiệu suất của quá trình truyền tải này là $\mathcal{H}_0$. Muốn tăng hiệu suất quá trình truyền tải lên 95\% nên trước khi truyền tải, nối hai cực của máy phát điện với cuộn sơ cấp của máy biến áp lí tưởng và cuộn thứ cấp nối với dây tải. Nhưng trong quá trình nối, do bị nhầm giữa cuộn sơ cấp và cuộn thứ cấp nên hiệu suất của quá trình truyền tải chỉ là 55\%. Giá trị của $\mathcal{H}_0$ là
	\choice
	{\True $\mathcal{H}_0 = 85 \%$}
	{$\mathcal{H}_0 = 88 \%$}
	{$\mathcal{H}_0 = 90 \%$}
	{$\mathcal{H}_0 = 95 \%$}
	\loigiai{
		\begin{longtable}{|p{0.17\textwidth}|p{0.17\textwidth}|p{0.17\textwidth}|p{0.17\textwidth}|p{0.17\textwidth}|}
			\hline %========================
			$P = h/s$ &
			$U_{SC}$ &
			$U_{TC}$ &
			$\Delta P$ &
			$P_t$\\
			\hline %========================
			100 &
			$U_1$ &
			$U_1$ &
			$100 - H$ &
			H\\
			\hline %========================
			100 &
			$U_1$ &
			$\dfrac{U_1}{k}$ &
			5 &
			95\\
			\hline %========================
			100 &
			$U_1$ &
			$kU_1$ &
			45 &
			55\\
			\hline %========================
		\end{longtable}
		Nhớ: $\Delta P\sim \dfrac{1}{U^2}\Rightarrow \dfrac{45}{5}=\left(\dfrac{U_1}{k.k.U_1}\right)^2\Rightarrow k=\dfrac{1}{\sqrt{3}}\Rightarrow \dfrac{100-H}{5}=\left(\dfrac{U_1}{k.U_1}\right)^2=\dfrac{1}{k^2}=3\Rightarrow H=85\%$.
	}
\end{vidu}
\begin{vidu}%[3P3-15.3K][Loai1] 
	Điện năng từ một trạm phát điện được đưa đến một khu tái định cư bằng đường dây truyền tải một pha. Cho biết, nếu điện áp tại đầu truyền đi tăng từ U lên 2U thì số hộ dân được trạm cung cấp đủ điện năng tăng từ 42 lên 177. Cho rằng chỉ tính đến hao phí trên đường dây, công suất tiêu thụ điện của các hộ dân đều như nhau, công suất của trạm phát không đổi và hệ số công suất trong các trường hợp đều bằng nhau. Nếu điện áp truyền đi là 3U thì trạm phát điện này cung cấp đủ điện năng cho
	\choice
	{214 hộ dân}
	{200 hộ dân}
	{\True 202 hộ dân}
	{192 hộ dân}
	\loigiai{
		Lập bảng và tóm tắt đề bài:
		\begin{longtable}{|p{0.22\textwidth}|p{0.22\textwidth}|p{0.22\textwidth}|p{0.22\textwidth}|}
			\hline %========================
			$P = h/s$ &
			$U$ &
			$\Delta P$ &
			$$P_t$$\\
			\hline %========================
			$36a + 42$ &
			$U$ &
			$36a$ &
			$42$\\
			\hline %========================
			$9a + 177$ &
			$2U$ &
			$9a$ &
			$177$\\
			\hline %========================
			$222$ &
			$3U$ &
			$4a = 20$ &
			$222 – 20 = 202$\\
			\hline %========================
		\end{longtable}		
		Dựa vào nhận xét: vì $P = const$ nên: 
		\begin{align*}
		36a + 42 = 9a + 177 \Rightarrow a = 5 \Rightarrow P = 222
		\xrightarrow{P=const}\Delta P\sim \dfrac{1}{U^2} \Rightarrow \dfrac{\Delta P_1}{\Delta P_2}=\left(\dfrac{U_2}{U_1}\right)^2.
		\end{align*}
	}
\end{vidu}
\begin{vidu}%[3P3-15.3K][Loai1] 
	\nguon{QG - 2018} Điện năng được truyền từ một nhà máy phát điện gồm 8 tổ máy đến nơi tiêu thụ bằng đường dây tải điện một pha. Giờ cao điểm cần cả 8 tổ máy hoạt động, hiệu suất truyền tải đạt $70\%$. Coi điện áp hiệu dụng ở nhà máy không đổi, hệ số công suất của mạch điện bằng 1, công suất phát điện của các tổ máy khi hoạt động là không đổi và như nhau. Khi công suất tiêu thụ điện ở nơi tiêu thụ giảm còn $72,5\%$ so với giờ cao điểm thì cần bao nhiêu tổ máy hoạt động?
	\choice
	{\True 5}
	{6}
	{4}
	{7}
	\loigiai{
		Chuẩn hóa: $P_0=100 $.
		\begin{longtable}{|p{0.22\textwidth}|p{0.22\textwidth}|p{0.22\textwidth}|p{0.22\textwidth}|}
			\hline %========================
			P &
			U &
			$\Delta P$ &
			$P_{t}$\\
			\hline %========================
			800 &
			U &
			240 &
			560\\
			\hline %========================
			100n &
			U &
			$3,75n^{2}$ &
			464,8\\
			\hline %========================
		\end{longtable}
		\begin{align*}
		\Delta P=\dfrac{P^2.R}{U^2}\Rightarrow \dfrac{\Delta P_2}{\Delta P_1}=\dfrac{P_2^2}{P_1^2}\Leftrightarrow \Delta P_2=240.\dfrac{100^2}{800^2}.n^2=3,75n^2
		\end{align*}
		Vậy: $P_2=\Delta P_2+P_{t2}\Leftrightarrow 100n=3,75n^2+406$
		$\Rightarrow n=5$ hoặc $n=21,7$.
	}
\end{vidu}
\begin{vidu}%[3P3-15.3K][Loai1] 
	\nguon{MH 2019} Điện năng được truyền tải từ nhà máy phát điện đến nơi tiêu thụ bằng đường dây tải điện một pha. Để giảm hao phí trên đường dây người ta tăng điện áp ở nơi truyền đi bằng máy tăng áp lí tưởng có tỉ số giữa số vòng dây của cuộn thứ cấp và số vòng dây của cuộn sơ cấp là $k$. Biết công suất của nhà máy điện không đổi, điện áp hiệu dụng giữa hai đầu cuộn sơ cấp là không đổi. Hệ số công suất của mạch điện bằng $1 $. Khi $k=10$ thì công suất hao phí trên đường dây bằng $10\%$ công suất ở nơi tiêu thụ. Để công suất hao phí bằng $5\%$ công suất ở nơi tiêu thụ thì $k$ phải có giá trị là
	\choice
	{$19,1$}
	{\True $13,8$}
	{$15,0$}
	{$5,0$}
	\loigiai{
		Ta có bảng số liệu sau
		\begin{longtable}{|p{0.17\textwidth}|p{0.17\textwidth}|p{0.17\textwidth}|p{0.17\textwidth}|p{0.17\textwidth}|}
			\hline %========================
			d &
			$P$ &
			$U$ &
			$\Delta P$ &
			$P_{tt}$\\
			\hline %========================
			Ban đầu  &
			110 (3) &
			10U &
			10 (2) &
			100 (1)\\
			\hline %========================
			Lúc sau  &
			110 (4) &
			kU &
			$\dfrac{10}{{\left({\dfrac{k}{10}}\right)}^{2}}$ (5) &
			110 - $\dfrac{10}{{\left({\dfrac{k}{10}}\right)}^{2}}$ (6)\\
			\hline %========================
		\end{longtable}
		%		\begin{itemize}
		%			\item Số liệu trong bảng sau được điền theo từng bước sau
		%		\item Ban đầu
		%		\begin{itemize}
		%		\item $(1)$ Ban đầu cho $P_{tt} = 100$.
		%		\item $(2)$ Suy ra $\Delta P=10 $.
		%		\item $(3)\ P=\Delta P+P_{tt}=110 $.
		%	\end{itemize}
		%		\item Lúc sau
		%		\begin{itemize}
		%		\item $(4)$ Vì P không đổi nên $P=110 $.
		%		\item Mà $\begin{cases}
		%		P=UI\\
		%		\Delta P=I^2R\\
		%		\end{cases}$ vì P không đổi, mà U tăng lên $\dfrac{k}{10}$ lần suy ra $I$ giảm đi $\dfrac{k}{10}$ lần, suy ra $\Delta P$ giảm $\left(\dfrac{k}{10}\right)^2$ lần nên
		%		\item $(5)\quad \Delta P =\dfrac{10}{\left(\dfrac{k}{10}\right)^2}$
		%		\item $(6)\quad P_{tt}=110-\dfrac{10}{\left(\dfrac{k}{10}\right)^2}$
		%	\end{itemize}
		%		\item Mà $\Delta P=0,05P_{tt}\Rightarrow \dfrac{10}{\left(\dfrac{k}{10}\right)^2}=0,05\left(110-\dfrac{10}{\left(\dfrac{k}{10}\right)^2}\right)$ suy ra $k=13,8$.
		%	\end{itemize}
	}
\end{vidu}
\begin{vidu}%[3P3-15.3K][Loai1] 
	Từ một trạm điện, điện năng được truyền tải đến nơi tiêu thụ bằng đường dây tải điện một pha. Biết công suất truyền tải điện năng đến nơi tiêu thụ luôn không đổi, điện áp và cường độ dòng điện luôn cùng pha. Ban đầu độ giảm điện áp trên đường dây bằng $20\%$ điện áp hiệu dụng ở trạm điện. Để công suất hao phí trên đường dây giảm $25$ lần so với ban đầu thì điện áp hiệu dụng ở trạm điện tăng bao nhiêu lần?
	\choice
	{\True $4,04$ lần}
	{$5,04$ lần}
	{$6,04$ lần}
	{$7,04$ lần}
	\loigiai{
		$\Delta P=\dfrac{P^2R}{U^2{\cos}^2\varphi};\quad P'=P-\Delta P\Rightarrow \begin{cases}
		P=UI;\ \Delta P=\Delta U.I\\
		\Delta U=0,2U_2\Leftrightarrow \Delta P=0,2P\\
		\end{cases}$
		\begin{longtable}{|p{0.22\textwidth}|p{0.22\textwidth}|p{0.22\textwidth}|p{0.22\textwidth}|}
			\hline %========================
			$P$ &
			$U$ &
			$\Delta P$ &
			$P'=const$
			\\
			\hline %========================
			$100$ &
			$U_1$ &
			$20$ &
			$80$
			\\
			\hline %========================
			$x$ &
			$U_2$ &
			$0,8$ &
			$x-0,8$
			\\
			\hline %========================
		\end{longtable}
		$80=x-0,8\Rightarrow x=80,8$
		Vì $R,\cos \varphi $ không đổi nên $\dfrac{\Delta P_1}{\Delta P_2}=\dfrac{P_1^2}{P_2^2}.\dfrac{U_2^2}{U_1^2}\Leftrightarrow \dfrac{20}{0,8}=\dfrac{100^2}{80,8^2}.\dfrac{U_2^2}{U_1^2}\Rightarrow \dfrac{U_2}{U_1}=4,04$.
	}
\end{vidu}